\PrelimChapter{Cronograma}{\topBanner}\label{cap:cronograma}
\vspace*{0.7cm}

El cronograma presentado en la \autoref{fig:cronograma} muestra la planificación estratégica definida para el desarrollo del proyecto. Este se estructuró en siete fases, ejecutadas de manera secuencial, con excepción de la séptima fase, que se desarrolló de forma transversal durante gran parte de la ejecución del proyecto. Las actividades se distribuyeron a lo largo de 24 semanas (seis meses).

La primera fase, denominada \textbf{Identificación de necesidades}, se llevó a cabo durante el primer mes. En esta etapa se identificaron las \ac{NPO} relacionadas con la autenticación y la gestión de identidades, constituyendo la base para el desarrollo de las fases posteriores.

Durante el segundo mes se ejecutó la \textbf{Fase 2: Caracterización de tecnologías}, en la cual se identificaron y analizaron diferentes tecnologías de servicios de directorio.

En el tercer mes, se desarrolló la \textbf{Fase 3: Selección de la solución}, donde se contrastaron las necesidades previamente identificadas con las tecnologías caracterizadas, con el propósito de seleccionar la alternativa más adecuada.

Durante el cuarto mes se llevó a cabo la \textbf{Fase 4: Diseño de la solución}, considerando tanto los requisitos identificados como la tecnología seleccionada para definir la arquitectura y el diseño de la propuesta.

En el quinto mes, se ejecutó la \textbf{Fase 5: Implementación del prototipo}, que consistió en el despliegue y configuración de un prototipo funcional en un entorno controlado.

Posteriormente, en el sexto mes se desarrolló la \textbf{Fase 6: Validación del prototipo}, en la cual se realizaron diversas pruebas para verificar el cumplimiento de los requisitos establecidos y evaluar el desempeño de la solución implementada.

Finalmente, la \textbf{Fase 7: Difusión y documentación}, se desarrolló desde la mitad del segundo mes hasta la finalización del proyecto. Esta fase permitió documentar de manera continua los avances, resultados y hallazgos obtenidos durante el desarrollo. Adicionalmente, se realizaron actividades de socialización en las que se presentó el progreso del proyecto ante la comunidad estudiantil, fomentando la retroalimentación y la divulgación de los resultados parciales alcanzados.

\begin{figure}[H]
	\centering
	\includegraphics[width=0.9\textwidth]{\figpathG cronograma.png}
	\caption{Cronograma de actividades del proyecto.}
	\label{fig:cronograma}
\end{figure}

